% =============================================================================
% =============================================================================
% HEADER										
% =============================================================================
% =============================================================================

\documentclass[11pt]{article}
\author{Sami Farhat}
\title{STE}

% -----------------------------------------------------------------------------
% Packages

\usepackage{titlesec} 				% to change the section font-size
\usepackage[margin=0.8in]{geometry} 	% for Margins
\usepackage{geometry}
\usepackage{hyperref}				% for links
\usepackage{multicol}				% for embedding lists in multiple-columns
\usepackage{enumitem}				% customize lists
\usepackage{fontspec}				% customise font used
\usepackage{tabularx}
\usepackage{soul}
\usepackage{color} 					% colouring text
\usepackage{array}					% bold table column
\usepackage{booktabs}				% http://ctan.org/pkg/booktabs
\usepackage{etoolbox}
\usepackage{float}
\usepackage[table]{colortbl}			% http://ctan.org/pkg/xcolor
\usepackage{tikz}					% for hierarchy structure
\usepackage[export]{adjustbox}
\usepackage{wrapfig}
\usepackage{graphicx}
\usepackage{longtable}
\usepackage{listings}
\usepackage{textcomp}
\usetikzlibrary{shapes,arrows}
\usepackage{subcaption}
\usepackage[perpage, bottom]{footmisc}		% restart footnote count on per page basis
\usepackage{calc}	% for math computations
\usepackage[nomessages]{fp} % http://ctan.org/pkg/fp
\usepackage{enumitem,amssymb}
\usepackage[fleqn]{amsmath}
\usepackage{mathtools}
\usepackage{unicode-math}
\usepackage{xcolor,listings}
\usepackage{textcomp}
\setmainfont{Helvetica}


% Graphics
% ---------------------
\graphicspath{ {res/} }

% Sections
% ---------------------
\setcounter{section}{0}
% no numbering of sections 
%\setcounter{secnumdepth}{0}

% Tables
% ---------------------
% increase margins of tables (1 is default)
\def\arraystretch{1.5}

\newcommand{\tblitembegin} {
\vspace{-\topsep}
\begin{itemize}[noitemsep, topsep=0pt,itemsep=-1ex,partopsep=1ex, parsep=1ex,leftmargin=*,]
}
\newcommand{\tblitemend} {
\end{itemize}
\vspace{-\topsep}
}
% for line breaks in table
\newcolumntype{b}{X}
\newcolumntype{s}{>{\hsize=.5\hsize}l}

\newcommand{\specialcell}[2][l]{%
\begin{tabular}[#1]{@{}l@{}}#2\end{tabular}}

\newcommand{\tblheader}{\textbf}
\definecolor{listinggray}{gray}{0.9}
\definecolor{lbcolor}{rgb}{0.9,0.9,0.9}

% Paragraphs
% --------
% remove paragraph indent
\setlength{\parindent}{0pt} 
\parskip=10pt

% Code 
% ----
\lstset{
language=Java,
upquote=true, 
%breakatwhitespace=false,         % sets if automatic breaks should only happen at whitespace
breaklines=true,                 % sets automatic line breaking
showspaces=false,
showstringspaces=false,
basicstyle=\ttfamily,
%numbers=left,
%numberstyle=\tiny,
frame=single,
%commentstyle=\color{gray}
keywordstyle=\color[rgb]{0,0,1},
commentstyle=\color[rgb]{0.133,0.545,0.133},
stringstyle=\color[rgb]{0.627,0.126,0.941}
}

% =============================================================================
% =============================================================================
%  CONTENT									
% =============================================================================
% =============================================================================

\begin{document}

% ============================================================================ 
% TITLE PAGE
% ============================================================================ 

\clearpage
%\maketitle
\thispagestyle{empty} % remove page number
\pagebreak 
\begin{titlepage}
\begin{center}
\vspace*{1cm}
{\Huge Software Testing} \\
\vspace{0.5cm}
{\LARGE STE\\}
\vspace{0.5cm}
\vspace{3.5cm}
{\LARGE Sami Farhat\\}
\vspace{0.1cm}
{\Large \# 1065452\\}
\vspace{0.3cm}
\vfill
\vspace{0.8cm}
Software Engineering Programme\\        
Computer Science\\
\vspace{0.5cm}
University of Oxford\\
United Kingdom\\
\vspace{1.0cm}         

\includegraphics[width=0.15\textwidth]{oxford-logo.png}
\end{center}
\end{titlepage}

% ============================================================================ 
% TABLE OF CONTENTS
% ============================================================================ 

\clearpage
\tableofcontents 
\thispagestyle{empty} % remove page number
\pagebreak

% ============================================================================ 
% CONTENT
% ============================================================================ 

\pagenumbering{arabic} % start numbering  



% TODO 
% ----
%
% Place the requirements in a table and number them to allow for clear referencing. 


% REQUIREMENTS 
\section{Review: System Requirements}


% Requests for improvements on the requirements. 
% TODO: would be great if we could write the requirements as we understood them.
% TODO: Talk about our impression from reading the system requirements and all the vaguenesses found in there. 

% what type of requirements document is this, are they informal? 
It is not evident at the outset what type of requirements document these requirements belong in. If they are meant to be informal, we would kindly ask for a formal requirements specification to guide our testing, as without one, there will be much room for vagueness and assumptions to be made. This could lead to testing being done on a false understanding of the requirements and this in turn would be costly both in terms of money and time. 	
\par

% Requirements are missing "definitions" section 
% ----------------------------------------------
The requirements are missing key sections that provide the completeness required by an SRS document to properly define the requirements of the project. 	
\par

For instance, we note the absence of a definitions section which defines key terms used in the requirements. It would be useful to point out that terms "will", "should", "must", and "may" conform to a predefined standard definition (e.g. IEEE Std 610.12-1990) to highlight the priority of said requirements and avoid ambiguity in their interpretation.
More significantly though, terms more specific to the application should be explicitly defined: 

\begin{itemize} 
	\item Client and customer are used interchangeably. We assume through contextual understanding that they refer to the same entity. However, it would be good to put that in writing. Again "Definitions" section before the requirements would be the place to have that (IEEE).
	\item System and application are used interchangeably: the 2 requirements before last state "They system must [..]" while others refer to "the application". Again, we assume through contextual understanding that the terms refer to the software being tested, but request the requirements writers specify this. 
\end{itemize}
\par

% vocabulary not defined before the requirements. 
% The requirements seem to be missing a section that defines key terms used in the the different requirements such as "will", "should", "must", and "may". This is common practice in SRS documents are defined prior to the listing of the requirements in order to more precisely depict the requirement priorities.
% It is good that the present requirements limit the key terms used to a subset of well-known ones (will, may, should, must) are used. This reduces the risk of vagueness emanating from (english understanding). 
% Some SRS documents will capitalise SHOULD, MUST, MAY to leave no room for confusion. 

% -------------------------------------------------
% Requirements organisation (bullets, no numbering)
% -------------------------------------------------
\subsection{Requirements organisation}
The requirements are bulletized, a numbering scheme would be much more apt as it allows testers and reviewers of the document the ability to reference individual requirements using their identifier. In some organisation, test methods or classes are prefixed with the requirement reference they address. 

We will take it on ourselves to label the requirements for our own workflow, it is suggested that any additions or alterations needed on the requirements be done on the table (reference to table here).

% --------------------------------------------------
% Requirements doc lacks of coherence 
% which lends to an overall lack of professionalism
% --------------------------------------------------
\subsection{Requirements coherence}

The first bullet in the list is not an explicit requirement, but an introductory sentence. This is misleading to the reader. 

Seventh bullet "The system must include a simple ‘help’ [...]" does not terminate with a full stop while the rest do. 
When outlining marital status taxing, we can see "Calculation of tax for Married people", "Calculation of tax for Single people", but the word 'divorced' is not capitalised, typo or lack of respect for divorced?  

% non-functional requirements. 
The only non-functional requirement related to the application's robustness. Should we assume there is no interest in security, performance? Surely this isn't the case for a Tax Calculation application. What platforms must the system run on? Must the language to be used be Java, or can developer preference overtake in the matter? 
\par

% * Non-functional requirements have not been detailed: do we expect the system to run on any platform? are there any speed / performance requirements that need to be catered for? For example, the requirements don't mention how many clients can be persisted on the app, which may imply (unlimited), but in such a case, it would be useful to reconsider the array design, as this does not provide enough flexibility going forward. 

% * Are there no security requirements? Perhaps the ability to allow locking the app with an app is desirable? After all, the software does hold people's tax codes, rates...


% --------------------------------------
% DEFECTS and VAGUENESS in requirements 
% --------------------------------------

\subsection{Defects and Vagueness}

In this section we address imprecise requirements and what we believe to be defects in the requirements. The provided requirements have been labelled (see Appendix \ref{sec:labelled-requirements}) to improve readability and accessibility of the report. 

% -----------------------------------------------------------------------------
% [1]
\subsubsection{R1 - The application must enable the storage of client’s personal details including. }

There is a lack of specificity on what storage means here. Does it refer to persistent storage? In which, case a form of database on disk would be required, or is it sufficient for it to be stored in memory during application runtime? The latter would not offer persistent for if the application or machine crashed.  

If persistent storage is to be used, it would be useful for the requirements to ouline the requirements on such database, are there non-function requirements such as performance metrics that the database should respect? 

% -----------------------------------------------------------------------------
% [2]
\subsubsection{R2 - The software should issue each customer a numeric identifier [...]} % which must be unique to that client and never be used by a different client (that is, if a client is deleted from the system then their ID should not be used again but has to die with them). }

This requirement is relatively more detailed than some of the others, given it is accompanied by an example detailing a scenario and how the application should behave with respect to that scenario. 

While 'die' is quoted, it is still an aggressive term to use in a requirement. 'Expire' is a reasonable alternative term, but we won't pedantize on this as there are more serious defects and vaguenesses in the requirements.

% -----------------------------------------------------------------------------
% [3]
\subsubsection{R3 - The PTMA should enable customer details to be updated.}

This requirement can be further expanded on, especially considering that it relates to significant chunks of the code. 
What details can a user update? Should it be allowed to update a customer's account number (ID)?
Is it required that there be a separate menu command for updating each of the following: first name, last name, tax code ? What about salary? Was this a desirable or required feature?    

% requiremrents lacking more explanation (e.g. 3)
% no expanding on how the update process should be. It seems that we want to allow the flexibility to update either of the customer's fields directly with its own command, but this does not link back to an evident requirement from the bulletised list. 

% -----------------------------------------------------------------------------
% [4]
\subsubsection{R4 - The application should be able to print lists of customer information [...] } % with users being able to select blocks of customers to be printed based on their unique identifiers.}

This requirements weakness only becomes evident when the application is tested and it becomes apparent that the Unique Identifiers generated for the created customers are never presented to the the application user for them to take note of. 

This requirement can further do with specificity to what type of block can be selected in the list to print. 
Should it be possible to use a comma-separated list of blocks to print out? This approach is not uncommon in other computer applications (e.g. specifying the list of pages to print on modern OSes is possible with "2-10,15-35"). This requirement could do with more specificity to the pages format. 

% -----------------------------------------------------------------------------
% [5]
\subsubsection{R5 - The application will calculate the tax amount to be paid by each client [...]} % and display this amount along with the other customer details. The amount of the tax is based on their taxcode and their current salary according to the rules below.}

This requirement is reasonably understandable. It is not however clear why it was separeted from the tax rules, seems like they could have been structured to follow each other. 

% -----------------------------------------------------------------------------
% [6]
\subsubsection{R6 - The system must include a simple 'help' system that lists all commands}
Assuming system refers to application, there's no ambiguity in the functionality requested in this requirement. 

% -----------------------------------------------------------------------------
% [7]
\subsubsection{R7 - The system must not crash if the user enters something that they are not meant to.}
Assuming system refers to the application, the requirement is clear. This requirement relates to the application's robustness and should, in a more formal document, be placed in the non-functional requirements section along with other requirements relating to performance, robustness, security. 

% -----------------------------------------------------------------------------
% [8]
\subsubsection{R8 - TaxRules}
[ Thus, for clarity, both of the following tax codes are valid. ] 
It would be extra helpful to provide a tax code that should not be accepted (e.g. 1080MT90). In fairness the provided examples were appreciated as they helped clarify the requriement.

% The alphabetic portion of the code may consist of multiple letters drawn from the
% following set, with their corresponding meanings:
The requirement does not explicitly elicit the fact that the different letters pertaning to marital status are mutually exclusive. Similarly, it does not do this for education and number of children. 

The requirements do not specify there must be a letter from the marital status in the tax code. Intuitively, an individual must be in one of those sections hence the tax code MUST have a letter from the marital status. This is not specified in the requirement but was assumed during testing. 

The requirement does not specify the application's behaviour for the case when duplicates of the same letter are found in the tax code. We would expect the tax code parser to complain and specify this as invalid, however others may interpret "DD1000" to be apply the Divorced taxing twice. 

% common to all next sections: 

The statement: "Married people with a current salary less than £12,000 will pay a tax rate of 90\% of their base tax amount.""

alone, implies that someone earning £4,000 must pay a rate of 90\%. This is not true, in the case of this application given the other statement: "Married people with a current salary of less than £5000 will pay a tax rate of 70\% of their base tax amount.". 

We assume the intention is to say: 

\begin{itemize}
	\item People earning between [0-5000[ should pay 70\% tax rate 
	\item People earning between [5000-12000] should pay 90\%
\end{itemize}

The requirements are not as well expressed here. 

% 8 - Married section 
The requirements for taxing Married people do not cover those who earn exactly £27,000. The third bullet ("less than") implies they are excluded from the 120\% tax rate, and the bullet right after ("above than") excludes them as well. 

% 8 - Single section 
Third bullet possibly mistyped. If it is not, then the third and fourth bullets are contradictory. We assume the intention was to say:  

\begin{itemize}
	\item People earning in bracket [11k-22k[ pay a rate of 125\%
	\item People earning 22k and above pay rate 132\%
\end{itemize}

If third bullet is adjusted to say "less than" instead of above, then this fixes the contradiction, but still leaves this requirement not covering the case where the taxee earns £22,000 exactly (same as Married section). 

% 8 - Divorced 
Same as single section. 

The following salaries / tax codes are not covered by the requirements: 
\begin{itemize}
	\item Married and earning £27,000
	\item Single and earning £22,000 (if other typo is fixed)* (see Single Section)
	\item Divorced and earning £24,000 (if other typo is fixed) *(see Divorced Section)
	\item Single Child and earning £10,400
	\item Two children and earning £9900
	\item More than 2 children and earning £9000
	\item FT education below 24yo and earning £3000
	\item FT education 24yo and above earning £3000
\end{itemize}

% ----------------------------------------------------------------------------------------------
% MISSING REQUIREMENTS ? 

% showing the ID of created clients 
Requirement 4 points to the need to allow users to see customers based on their IDs, but there is no reference to how the application-generated IDs are presented to the user. This was found to be an issue during exploratory testing, and was sourced back to both a defect in the implementation and absence in the requirements. 

% no requirement to delete clients. 
There is no requirement about the user's ability to delete customers from storage. This is tacitly mentioned as part of the second requirements relating to ID uniqueness. We believe it is then an oversight on the requirements writers. 

% Edge case
Is it implicity that negative salaries are not a thing? How does the code deal with that? 

% ----------------------------------------------------------------------------------------------
% Adjustments and assumptions we make on the requirements. 
\subsection{Assumptions}
\begin{itemize}
	\item We assume that earning less than zero (Paying?) person, pays no tax. i.e. min range = 0. 
\end{itemize}

% END OF REQUIREMENTS
% ----------------------------------------------------------------------------------------------

% Talk about our impression from reading the system requirements and all the vaguenesses found in there. 

\section{Code Review}

\subsection{Code Read}
% Detail what notes were taken when reading the code and what was noticed from afar
% We should say that reading the code is not great in that it usually focuses our testing on some things and push us not to test, things that 
% we expect to be correct (given we know the implementation). As such, it is always better to try and make no assumptions about the contents 
% of the code when writing the tests. 

% DEFECT: There's somewhat of an assumption being made that a users salary always fits in Integer. What if they earned more than 2^31 - 1 ? 

\subsection{Code Refactoring}

% md5_refactor_check
% ------------------
% Talk about how we want to prevent our code refactoring from actually fixing hidden bugs. Of course we cannot practically be sure that it does not, but we have done is write a script that compares the stdout from certain sequences of commands from both the original code and the modified-refactored code. The script generates the md5 hash from each of the outputs. They should match. If they don't this points to us having modified the current output. This script has been added as a post-commit hook, which means that on commit made a Jenkins build runs.  


\subsubsection{Changes in TaxEngine.java}

% [1] create the exception type TaxEngineException
[1] First off, we are displeased with the error reporting of this class, it does not make testing for error conditions any easier. We create an "exceptions" package and add TaxEngineException which extends RuntimeException in this package. We will expect this exception object to be thrown in TaxEngine methods where errors occur. 

% why choose RuntimeException
We chose RuntimeException as it then does not require all methods to add the "throws TaxEngineException" which will be added to a lot of sections of the code causing many other changes. We opted for the most gains for the least upsets to the code here. 

% [2] split out the different branches in taxAmount() to smaller more testable methods
% [2] The contents of every block similar to `if (strTaxCode.indexOf(Constants.MARRIED_CODE) >0 ) {` were moved to private methods, making for 8 in total (see below). As these methods are private, they cannot be tested without the usage of reflection or other Java dark arts. As such, we further create 3 public methods (see below).

In taxAmount(String, int), the changes involve replacing the contents of the different if blocks with the right method calls to taxBasedOnMaritalStatus(), taxBAsedOnChildren() and taxBasedOnEducationalStatus() each time passing the proper enum. 

% Just trying to explain that separating the big method into smaller ones that are individually testable is better, as it allows us to more quickly/easily locate bugs when the different tests fail. 
From a testing point of view, we still want to write tests on taxAmount with the different combinations of TaxCodes and input parameters, checking each time if the return is as expected. Doing  

But in addition to the above, we also write different tests for each of the [taxBasedOn...] method. With simpler methods in TaxEngine, and tests targeting those simple methods, we can more quickly identify the source of bugs when tests fail. The alternative 

Basically if there was something wrong with how we compute salaries for Married people, then I'd expect this to manifest in the test for taxAmount() and
in taxBasedOnMarriage(), these two methods failing together should indicate to the tester or developer checking the tests that it is likely that the logic
in Marriage taxing is causing that in taxAmount to fail. Having both higher level and lower level tests is useful in that way. 

% -----------------------------
% Methods added: 
% -----------------------------
* taxBasedOnMaritalStatus
* taxBasedOnChildren
* taxBasedOnEducationStatus  
% -----------------------------
* taxMarriedCode
* taxSingleCode
* taxDivorcedCode
* taxOneChildCode
* taxTwoChildrenCode
* taxMultipleChildrenCode
* taxFullTimeStudentUnder24
* taxFullTimeStudentOver24
% -----------------------------

\subsection{Changes in Main}

% this is not a change that we've made, just a note that we have redirected stdin and stdout of the application 
% in order to capture and analyse its output and to control its input. 
* Allow for the selection of PrintStream and UserInputStream to Main. If not set then System.in and System.out are used. 
We could have also just used System.setOut() and System.setIn(), but our approach should be even more flexible as it doesn't affect our usage of System.out in the testing framework. 

% * Make methods non-static. The reason for them being static is unclear at present; it could be that a behaviour similar to the one provided by a Singleton pattern design is the sought one. We decide - aginast our natural tendency - not to change the static types of the methods since we can still work with them being static. 

* Allow for Main's start/stop execution to be controlled via methods. Currently, to stop main, the user must press 'q' which then leads to System.exit(). If running a test suite, this will cause the whole test suite to exit. We replace System.exit() with logic that causes the method to return. 
 -	In a normal running App,   when the method returns, the app will reach the end of Main and exit naturally (and gracefully), when run under a test suite, the framework will continue to run after the method has returned which is what we wanted in the first place. 
  - Add method `stop()` and member field `running` which when set to false causes the User Input loop to exit, leading to stop of the program.

% print the ID of the created Customer otherwise there is no way to know how to list that specific customer. (see newCustomer, has the 'id' of the new customer printed.)

% add method getCustomer in Main to aid with testing. As we would prefer getting the customer object from main directly rather than going through the route of listing customers with that ID, and serializing a customer object from the printed fields to use in our tests. e.g. adding "Sami Farhat TB1000 10000", then the ID is output, the we would have to, use 'l' <id>, and get the output from the program and serialise a Customer object from the output. Also, a bug in the listing would cause all tests depending on getting a customer from the database to fail. To separate the tests better, we create a method in Main that returns a customer object directly from the data structure in AllCustomers given a CustomerID. 


% Change in AllCustomers.java

Pass PrintStream 'out' via DI injection too. If none provided, then System.out is used (default). This is needed since the AllCustomers class is responsible for printing lists of customers.  


\section{Test implementation}


% TODO:  use surefire report
\subsection{Unit testing}

* we use display names to give more readable output for tests
* we use parametrized tests to allow for high flexibility in generating the testCases. --> Explain the alternative where we would have to write an explicit test case for every input type. And further explains, why running a loop with asserts() in there is not equivalent to our current approach. (Basically the answer to the question, why did we use Paramterized testing). 
* we acknowledge that some regex guru may come up with a regex that replaces the 3 we currently use right now. That's not a problem, we opted for readability anyway ;) 
* we've made some effort to make the tests generalised in that. If some of the requirements change (change the letter used for a symbol) or if more letters are added, then we wouldn't have to overhaul the whole Test class but merely make some select changes. Of course, some tests take @ValueSource which are hardcoded strings to test. In case of change in requirements, those especially should be carefully reviewed (one by one). 


\subsection{Cucumber testing}

% not A LOT OF REFACTORING was done 
Explain that, we questioned how much refactoring we should be doing before starting the testing and opted on the conservative end of the spectrum deciding not to do too much and fix the existing code. Our role is to test it, and it may even be that we have been hired to highlight how many defects existed so that administration takes action on the developers (oops), in which case we must stick to our defined role and do testing the best we can given what we are given. 

Other changes and refactoring could have been done to Main, such as setting up an event communication loop that would allow the thread in which Main ran to synchronize with the cucumber thread that is doing the testing. We opted against making major changes. 

Method getCustomer was added to main, this was needed for testing. It acts as a proxy to allCustomers. An alternative approach could have been to return a getter to the allCustomers private field, or to simply make that field public. But as only one feature of allCustomers is needed for testing (at the time of writing at least), we went ahead an exposed the single method we care for from allCustomers. 

Many of our Scenarios (if not all) depend on applicationHasStarted(). This step starts the "Main" application. By design "Main" waits for user input to proceed and will block the execution thread until exits. To stop this, we make applicationHasStarted start "Main.run()" in a separate thread. 

This introduces some timing challenges. Subsequent steps in our framework need to work with Main by injecting user input and monitoring the output, but the steps need to know if "Main" is currently accepting of user input and if it has printed the output it requires. 

One solution to this is to add a boolean field member `waitingForCommand` in "Main" which indicates that the app is currently waiting for the next command to execute. The framework can then know for certain if it can safely inject new input to System.in and whether it can assume that Main has printed all it needed to already. This latter is ensured by the fact that Main will not have anything more to print() when it is blocked waiting for a new command from user input. 

% note on checking uniqueness
It's near impossible to fully confirm that the application is indeed returning unique identifier without incurring high computational costs. For example, the application may be coded to return a Random int ID in the range [1-100] for each customer. A test that creates 5 customers would still only have 10\% chance of hitting a uniqueness case: 
$1 - P(100,5) = 0.10$. This is known as the Birthday problem % https://en.wikipedia.org/wiki/Birthday_problem]. 
An application returning a random Integer. 

% The concurrency issues observed 
Whilst doing cucumber testing we noted, that the Cucumber suite ran tests in parallel and while this is desired to reduce the time it takes to run tests, it introduced multiple threading issues. With Main class being static, all threads effectively shared the same PrintStream  - as well. [ More to say ? ]

We then decided to refactor both our test suite, which at the time comprised of one class containing all Step Definitions, and the Main class. Main was moved from being static based, to instance based; all static modifiers were removed, a method run() was created. Main constructors were introduced: a default one that would use the standard (default) stdin and stdout, and another that can have those provided as a dependency injection. This greatly increased flexibility and our ability to test the different Main instances independently. 

% NOTE: refactoring Main to become non-static introduces an issue with the customer identifier index being static 

% Refactoring of our Cucumber test suite. 
We moved from having a single StepDefs.java class that contained all steps to creating a separated FeatureXStepDefs class for each feature undergoing acceptance testing. Those classes shared a lot in common in terms of the helper functions they would need to call, but also of the member variables they would use to test the app's behaviour. For this reason, we moved these common members and helper functions into their own class BaseStepDefs, which then became the inheritance parent of each FeatureXStepDefs. 

% challenge with ApplicationHasStarted 
As all features defined in .features rely on the Given "Application has started", we first thought it logical to move the corresponding step definition method into BaseStepDefs; alas, cucumber complained that a class defining a Step Definition cannot be inherited from, so this could not be done. Our workaround, was to defined applicationHasStarted() as a helper in BaseSteps and have a separate regex "application has started" for each feature:
% * in 1_ClientAdd.feature, there's  Application has started1"
% * in 2_UniqueID.feature, there's Application has started2"
... 
Each of those step definitions then calls the method applicationHasStarted() from BaseStepDefs, which avoids code redundancy. 

This approach may not be perfect, but (1) it does the job, (2) it avoids code duplication, (3) it is not very confusing. 

\subsection{Testing calculated tax values}

Time was taken to decide how to proceed with verifying correctness of tax values. JUnit offers @ParametrizedTest annotation that allows the calling of a test function multiple times with different parameters, each time checking the correctness of the output given the input parameters. This is not only useful but needed, it is far too costly time and effort wise to write a new tests for each combination of tax code and salary. (Provide calculation estimate here on how many there are approximately).
If the tests were to be spelled out individually (one test for each taxcode-salary combination), then it would require even more testing time if a new taxcode 'letter', or a new salary bracket were to be added in the future. (e.g. adding a new taxcode letter for people with 3 children exactly, would require xxxx tests more). 

For the above reason, we looked into parametrized tests using Cucumber and found this to be possible using the "Examples" keyword. 
An "Examples" section placed just below a scenario definition determines the combinations of parameters to be used when running this Scenario, for our tax verification purposes we use the table below: 

firstName	| lastName	| taxCode 	| Salary 	| expectedTax
Tobias 		| Mann 		| TB1000 	| 10000 	| ....
Sophie		| Caseby	| S2000 	| 20000		| ....

This entails the test will be run with Tobias as firstName, Mann as lastName... for the first run of the scenario. Then for the second run "Sophie" will be used for the firstName and so on. 
This approach allows us to define fewer scenarios that can take in more combinations increasing our scenarios' readability all the while reducing the time needed to test the large number of combinations we need to. 

The changes needed for this to work on the Scenario's side of things is to use <firstname> and <lastname> in the steps in the feature file instead of hardcoding the values, and to add String parameters to the Java StepDefinitions.

% ExamplesTableGenerator 
% 
Explain how the tool is provided to help generate the Examples table that can be plugged into the .feature file. And how after some careful reading, it was determined that it's not a great solution (also it is not implemented in cucumber) to have the feature file load a separate examples file, because we want to have it all be contained. 

While it may be argued that the approach taken here is overkill or not in the sprits of Behaviour Driven Testing, we would argue that it is our role to detect whatever defects we can and that the implications of defective code are very costly in our case, we have been asked to find the most defects using cucumber which is a BDD tool and we have taken the necessary approach that we believe achieves the objective. 


\subsection{List of defects}

% Requirements defects 
% Logic defects
% User Interface defects
% Code style defects


% ============================================================================ 
% APPENDIX
% ============================================================================ 

\appendix

\setcounter{section}{0}
\pagebreak 
\clearpage
\thispagestyle{empty} % remove page number
\vspace*{9cm}
\begin{center}
{\bf \LARGE Appendix}
\end{center}
\vfill
\pagebreak

% \section{More details}


% =============================================================================
% END DOCUMENT
% =============================================================================

\end{document}

% =============================================================================
% CHECKS
% =============================================================================

% [  ] Check spelling
% [  ] Check with "grammarly"

% =============================================================================
% DRAFT
% =============================================================================